% !TEX root = vkr.tex
\section{Обзор методов и инструментов}
\label{sec:review}

Обзор охватывает направления, необходимые для расширения PySATL-CPD:
офлайн-обнаружение разладок, применение алгоритмов к последовательно
поступающим данным, оценку качества и подготовку воспроизводимых
экспериментов. Методы сравниваются по доступным данным, виду результата,
предположениям и требованиям к интеграции. Отдельное внимание уделяется
различию между существующими онлайн-обёртками и оконным запуском
офлайн-алгоритма. Рассмотренные решения служат основанием для выбора
интерфейсов и протокола проверки; реализация и экспериментальное
сравнение относятся к следующим этапам работы.

\subsection{Офлайн- и онлайн-постановки}

Пусть наблюдается последовательность $x_0,\ldots,x_{n-1}$.
Разладка означает изменение вероятностного механизма, порождающего
данные: например, среднего, дисперсии или другой характеристики
распределения. Отдельный выброс сам по себе не обязательно соответствует
переходу в новый режим. Выбор алгоритма поэтому зависит от того,
какие изменения требуется обнаруживать и какие предположения
допустимы для рассматриваемых данных.

В офлайн-постановке доступен весь анализируемый ряд или фрагмент,
а результатом служат оценки границ сегментов. В онлайн-постановке
при обработке наблюдения $x_t$ доступны только $x_0,\ldots,x_t$;
существенно не только положение изменения, но и время реакции.
Офлайн-сегментация представлена, в частности, в
\texttt{ruptures}~\cite{ruptures}, а примером последовательного
вероятностного подхода является BOCPD~\cite{adams2007}.

Для проектируемого расширения целесообразно явно хранить
$(\hat\tau,d)$: оценку позиции разладки и индекс наблюдения,
после обработки которого выдан сигнал. Принимается нулевая индексация,
а $\hat\tau$ обозначает первое наблюдение нового сегмента.
Начало и конец ряда не считаются обнаруженными изменениями.
Такое соглашение позволит сравнивать обычную сегментацию и её
оконную адаптацию, не смешивая ошибку локализации с задержкой.
Метод поиска одной границы также должен отличать оценку её позиции
от решения о наличии изменения: максимум статистики можно найти
и на однородном ряду.

\subsection{Подходы к обнаружению разладок}

\subsubsection{Сегментация по функции стоимости}

Распространённый офлайн-подход состоит в поиске разбиения,
минимизирующего сумму стоимостей сегментов со штрафом за число границ:
\[
 \min_{m,\,\tau_1<\cdots<\tau_m}
 \left\{\sum_{j=0}^{m}c(x_{\tau_j:\tau_{j+1}})+\beta m\right\},
 \qquad \tau_0=0,\quad\tau_{m+1}=n.
\]
Здесь $x_{a:b}$ --- наблюдения с индексами из полуинтервала $[a,b)$,
$c$ --- стоимость описания сегмента, $\beta$ --- штраф.
Функция стоимости задаёт модель изменений, а процедура поиска
определяет способ выбора границ. Это разделение позволяет менять
статистическую модель независимо от стратегии сегментации.

PELT использует динамическое программирование с отсечением
кандидатов~\cite{killick2012}. При выполнении условий отсечения
сохраняется минимум выбранного критерия. Ожидаемая линейная сложность,
указанная авторами, требует дополнительных условий и не является
безусловной гарантией. Результат зависит от стоимости и штрафа,
поэтому само название алгоритма не определяет чувствительность
ко всем видам разладок.

Бинарная сегментация последовательно ищет одну границу, разделяет
по ней фрагмент и повторяет поиск в полученных частях~\cite{binseg}.
Она позволяет использовать процедуру поиска одного изменения
для нескольких разладок, но раннее разбиение ограничивает дальнейшие
решения. Для PySATL-CPD это аргумент в пользу разделения базового
детектора и стратегии многократного поиска. PELT и бинарная
сегментация рассматриваются как ориентиры для сравнения интерфейсов;
реализация всех рассмотренных методов в рамках семестра не требуется.

\subsubsection{Ядерные и графовые подходы}

Методы сравнения распределений расширяют возможности моделей,
ориентированных только на изменение среднего. Например, функция
\texttt{CostRbf} в \texttt{ruptures} оценивает однородность сегмента
через среднее в пространстве признаков, заданном ядром~\cite{kernel}.
RBF-ядро позволяет учитывать более общие различия распределений.
При этом результат зависит от масштаба признаков и параметров ядра;
непараметрический подход также требует обоснованного выбора настроек.

Графовый метод Chen и Zhang~\cite{chen2015} строит граф сходства
наблюдений и оценивает разделение по числу рёбер между частями ряда.
Необычно малое число таких рёбер указывает на различие групп.
Метод не задаёт конкретное параметрическое семейство распределений,
но зависит от расстояния, графа и калибровки статистики.
Перестановочная калибровка требует учитывать предположения
об обменности; для временно зависимых данных обычная перестановка
не обеспечивает их автоматически. В данной работе графовый метод
выбран первым офлайн-алгоритмом для интеграции. Его роль ---
предоставить рабочую основу для проверки общего интерфейса,
оконной обёртки и метрик. Подробности конкретной реализации
будут определены на следующем этапе.

\subsubsection{Последовательные методы}

BOCPD рекурсивно обновляет распределение длины текущего участка
без изменений~\cite{adams2007}. Предиктивная модель наблюдений
сочетается с функцией риска, задающей вероятности смены режима.
В исходной постановке наблюдения внутри сегмента условно независимы,
а параметры разных сегментов независимы между собой.
Апостериорное распределение ещё не задаёт список тревог:
для экспериментального сравнения требуется отдельное правило решения.
Без усечения число рассматриваемых длин сегмента растёт со временем.

Последовательный метод обновляет состояние при поступлении данных,
тогда как оконная адаптация повторно анализирует ограниченный фрагмент.
Поэтому сравнение должно учитывать частоту решений, задержку,
ложные тревоги и стоимость обработки. Доступные онлайн-алгоритмы
PySATL-CPD могут служить базовыми линиями; состав небольшого набора
для экспериментов следует выбирать по совместимости моделей
с выбранными данными, а не только по наличию реализации.

\subsection{Оконная обработка и обобщение обёрток}

Оконная адаптация связывает офлайн-детектор с последовательным входом:
обёртка накапливает данные, формирует фрагменты и преобразует
полученные границы в события исходного ряда. Её поведение должно
определяться отдельно от статистики конкретного алгоритма.
Для окна длины $w$, заканчивающегося на индексе $t$, локальная
граница $k$ преобразуется по правилу
\[
 \hat\tau=t-w+1+k,\qquad d=t.
\]
Например, окно $[80,100)$ и локальная граница $12$ дают позицию $92$,
но результат становится доступен после наблюдения $99$.
Это пример соглашения об индексах, а не результат эксперимента.

В текущих обёртках PySATL-CPD \texttt{skipping} пропускает наблюдения
по условию, а \texttt{batching} сворачивает накопленный блок в одно
значение и передаёт его онлайн-алгоритму~\cite{pysatlwrappers}.
Их различия с проектируемым механизмом приведены
в таблице~\ref{tab:wrappers}.

\begin{table}[htbp]
\centering
\caption{Способы обработки последовательного входа}
\label{tab:wrappers}
\small
\renewcommand{\arraystretch}{1.2}
\begin{tabularx}{\textwidth}{@{}L{0.22\textwidth}YY@{}}
\toprule
Механизм & Действие с данными & Что получает алгоритм \\
\midrule
\texttt{skipping} & Отбор наблюдений по условию & Отдельные сохранённые наблюдения \\
\texttt{batching} & Накопление и свёртка неперекрывающегося блока & Одно агрегированное значение \\
Скользящее окно & Хранение фрагмента длины $w$, запуск с шагом $s$ & Полный фрагмент для офлайн-анализа \\
\bottomrule
\end{tabularx}
\end{table}

Алгоритм \texttt{ruptures.Window} решает другую, более конкретную
задачу: вычисляет выигрыш от разделения соседних частей окна
и затем выбирает пики этой величины~\cite{ruptureswindow}.
Само наличие скользящего окна в названии метода не означает,
что он является универсальным адаптером произвольного детектора.
Для PySATL-CPD требуется именно независимый механизм запуска.

Из сравнения следует, что обобщать целесообразно конфигурацию,
сброс состояния и правила композиции. Типы результатов при этом
могут различаться: скалярная статистика онлайн-алгоритма и набор
границ офлайн-метода имеют разные назначения. После пропуска
или агрегации наблюдений необходимо сохранять отображение
на исходные индексы; одного смещения окна уже недостаточно.

Перекрывающиеся окна могут повторно обнаруживать одну разладку.
Правило объединения должно использовать только уже доступные
результаты и сохранять момент первого сигнала. Уточнение позиции
позже не должно уменьшать ранее наблюдавшуюся задержку.
Также необходимо определить прогрев, неполное последнее окно,
сброс и поведение на коротких рядах. Эти случаи входят в контракт
и набор будущих модульных проверок.

При стоимости одного запуска $C(w)$ повторная обработка имеет
порядок $\lceil n/s\rceil C(w)$ без переиспользования вычислений.
Меньший шаг повышает частоту решений и нагрузку; увеличение окна
даёт больше наблюдений, но может объединять несколько режимов.
Память буфера ограничена длиной окна, а рабочая память детектора
может быть больше. Следовательно, $w$ и $s$ следует рассматривать
как параметры эксперимента наряду с настройками самого алгоритма.

\subsection{Метрики и правила оценки}

\subsubsection{Классификационные метрики}

Оценка найденных границ требует допускать ошибку по времени.
Если истинная граница равна $\tau$, предсказание $\hat\tau$
считается допустимым кандидатом на совпадение при
$|\hat\tau-\tau|\leq\delta$. Сам допуск $\delta$ должен быть
зафиксирован до итогового сравнения и выражен в понятных единицах.
В бенчмарке Van den Burg и Williams используются сопоставление
с допуском и учёт нескольких вариантов разметки~\cite{vandenburg2022}.

Для основного протокола предлагается взаимно однозначное сопоставление:
каждая истинная и найденная граница участвует не более чем в одной
паре. Сначала максимизируется число допустимых пар, затем среди
равноценных вариантов минимизируется суммарная ошибка локализации.
Это предлагаемое правило для проекта; его следует явно отличать
от режима точного воспроизведения конкретного бенчмарка.
Число пар задаёт $TP$, несопоставленные предсказания --- $FP$,
несопоставленные истинные границы --- $FN$. Тогда
\[
 P=\frac{TP}{TP+FP},\qquad
 R=\frac{TP}{TP+FN},\qquad
 F_1=\frac{2TP}{2TP+FP+FN}.
\]
Лишние сигналы около одной границы не должны увеличивать число
правильных обнаружений. Для рядов без разладок отдельно учитываются
ложные сигналы; случаи нулевых знаменателей требуют явно
заданного соглашения. В итогах полезно показывать показатели
по рядам и отдельно объединённые счётчики, поскольку длинные
последовательности иначе могут определять общий результат.

В PySATL-CPD уже имеются классификационные метрики и средства
сопоставления~\cite{pysatl,pysatlmatching}. Поэтому расширение
начинается с проверки их соглашений: допустимых отклонений,
повторных срабатываний и использования времён обнаружения.
Наличие функции ещё не означает её применимость к новому
представлению результата без адаптации.

При сильном дисбалансе точек с разладкой и без неё обычная
доля правильных ответов малоинформативна. Если метод выдаёт
оценки или допускает перебор порогов, дополнительно можно строить
PR-кривую. В библиотеке уже есть визуализатор PR-AUC~\cite{pysatlpr};
следует проверить способ построения кривой и обработки крайних случаев.
Average precision и площадь по правилу трапеций различаются,
что явно отмечает документация scikit-learn~\cite{averageprecision}.
Одного списка границ при единственном пороге для построения
полной PR-кривой недостаточно.

\subsubsection{Локализация и качество сегментации}

Precision и recall не различают близкую и предельно допустимую
ошибку внутри выбранного допуска. Поэтому для сопоставленных
границ предлагается дополнительно вычислять среднюю или медианную
абсолютную ошибку $|\hat\tau-\tau|$. Её необходимо публиковать
вместе с полнотой: метод, нашедший только простые изменения,
может иметь малую ошибку среди успешных обнаружений.

Для оценки всего разбиения подходит covering из
CPD-бенчмарка~\cite{vandenburg2022}: каждому истинному сегменту
сопоставляется предсказанный с максимальным отношением пересечения
к объединению, после чего значения усредняются с весами по длине
истинных сегментов. Этот показатель дополняет оценку событий,
учитывая качество покрытия ряда сегментами. При нескольких
аннотациях можно оценивать результат относительно каждой из них
и затем усреднять, сохраняя информацию о расхождении разметок.

Другие специализированные меры --- расстояние Хаусдорфа между
наборами границ~\cite{hausdorff} и индекс Рэнда для пар точек,
отнесённых к одинаковым или разным сегментам~\cite{randindex}.
Первая мера подчёркивает наихудшее расхождение, вторая оценивает
согласованность разбиений. Их полезно рассматривать как дополнительные:
для минимального протокола достаточно событийных метрик,
ошибки локализации и covering.

\subsubsection{Последовательное обнаружение и стоимость}

Для онлайн-методов и оконных адаптаций требуется оценивать
задержку $d-\tau$. При её расчёте учитываются только сигналы,
выданные после соответствующей истинной границы; преждевременные
тревоги требуют отдельного учёта. Ретроспективная близость
оценённой позиции не делает ранний сигнал правильным обнаружением.
Пропуски также должны оставаться видимыми рядом со средней задержкой.

На однородных рядах следует измерять число ложных тревог
на единицу длины и длину работы до первой ложной тревоги.
Если тревога до конца ряда не возникла, наблюдение цензурировано:
конец записи не следует выдавать за фактическое время срабатывания.
Порог отдельного офлайн-теста не гарантирует тот же уровень
ложных тревог при многократном применении к перекрывающимся окнам.

Итоговое сравнение должно включать время обработки и пиковую память
при фиксированном окружении. Для оконной адаптации полезны также
число запусков и средняя стоимость обработки одного наблюдения.
Таким образом, качество описывается несколькими группами показателей,
а преимущество метода по одной метрике не заменяет остальные.

\subsection{Наборы данных для проверки}

Основной ориентир --- Turing Change Point Dataset (TCPD),
содержащий ряды, сведения об их происхождении и ручные аннотации;
TCPDBench предоставляет инфраструктуру запуска и анализа
экспериментов~\cite{tcpd,tcpdbench}. Эти проекты выполняют
разные функции: набор данных задаёт объекты проверки,
а бенчмарк --- способ организации сравнения.

Для первого набора кандидатов подходят одномерные \texttt{nile}
и \texttt{well\_log}, а также многомерный
\texttt{bee\_waggle\_6}~\mbox{\cite{tcpdnile,tcpdwell,tcpdbee}}.
Первый содержит годовые значения стока Нила у Асуана,
второй --- геофизические измерения скважины,
третий --- координаты и направление движения пчелы.
Такой набор позволяет рассмотреть разные размерности и свойства
наблюдений. Окончательный выбор должен учитывать длину рядов,
число размеченных изменений, расстояние между ними и поддержку
типа данных выбранным алгоритмом. Угловой признак требует
отдельного решения о представлении и расстоянии.

В TCPD встречаются несколько разметок одного ряда; их нельзя
сводить к единственной бесспорной истине без описания процедуры.
Для сопоставимости с опубликованным бенчмарком потребуется отдельный
режим, повторяющий его правила агрегации precision и recall
по аннотаторам~\cite{vandenburg2022}. В основном протоколе
целесообразно сохранять исходные варианты разметки.
При загрузке также проверяются индексация, пропуски, происхождение
и условия использования каждого набора: лицензия кода репозитория
не заменяет условий для всех данных~\cite{tcpd}.

Последовательная подача записанного ряда позволяет проверить
причинность обработки, но не превращает данные в настоящий поток
с исходными задержками поступления. Дополнительно нужны
синтетические ряды с известными границами: без изменений,
со сменой среднего или дисперсии, с близкими разладками
и изменениями около края окна. Они позволят отдельно проверить
контракт и сложные случаи, не подменяя оценку на реальных данных.

\subsection{Инструменты и организация экспериментов}

В просмотренной версии PySATL-CPD офлайн-ядро ещё не содержит
реализации алгоритма. Онлайн-ядро предоставляет обработку
наблюдения, конфигурацию, состояние и сброс~\cite{pysatl}.
Поэтому первый этап интеграции должен определить результат
офлайн-детектора и его совместимость с данными библиотеки.
Библиотека \texttt{ruptures} полезна как внешний источник
сопоставимых реализаций и пример разделения стоимости и
стратегии поиска~\cite{ruptures}. TCPDBench показывает организацию
конфигураций, запусков и последующего анализа~\cite{tcpdbench}.

Для будущих экспериментов предлагается фиксировать версии данных
и кода, параметры предобработки, начальные состояния генераторов
случайных чисел,
настройки алгоритма и обёртки. Сохраняться должны как итоговые
показатели, так и найденные позиции и времена сигналов.
Это позволит пересчитать метрики без повторения всех запусков.
Параметры и пороги подбираются на выделенных данных;
итоговая оценка проводится отдельно. Предобработка при
последовательном запуске не должна использовать будущие наблюдения
или статистики, оценённые по всему тестовому ряду.

Сравнение офлайн-метода на полном ряде, его оконной адаптации
и ограниченного набора онлайн-методов должно явно учитывать
разный доступ к данным. Полный ряд служит ориентиром для оценки
сегментации, а задержка сравнивается только у последовательных
вариантов. Если базовый детектор ищет одну разладку на фрагменте,
это ограничение необходимо учитывать при выборе сценариев.
Проверки правильности индексов и метрик предшествуют измерению
качества; количественные результаты на обзорном этапе не заявляются.

\subsection{Анализ обзора и выводы для работы}

Рассмотренные подходы решают связанные, но разные задачи.
Оптимизационные методы определяют разбиение по функции стоимости,
непараметрические методы сравнивают распределения через выбранное
представление данных, последовательные методы обновляют решение
по мере поступления наблюдений. Оконная обработка позволяет
использовать офлайн-детектор при последовательном входе,
но требует самостоятельных правил выдачи событий и оценки задержки.
Следовательно, основная задача расширения библиотеки заключается
в согласовании этих компонентов и воспроизводимой проверке.

Таблица~\ref{tab:decisions} связывает результаты качественного
сравнения с задачами работы.

\begin{table}[htbp]
\centering
\caption{Выводы обзора и направления реализации}
\label{tab:decisions}
\small
\renewcommand{\arraystretch}{1.2}
\begin{tabularx}{\textwidth}{@{}L{0.26\textwidth}Y@{}}
\toprule
Направление & Вывод и ожидаемый результат \\
\midrule
Офлайн-детектор & Общий контракт результата; графовый метод как первая совместимая реализация. \\
Окна и обёртки & Независимый запуск на фрагментах; явные индексы, время сигнала, объединение повторов и сброс. \\
Метрики & Адаптация имеющихся средств; взаимно однозначное сопоставление, covering, локализация, задержка и ложные тревоги. \\
Данные & Небольшой набор TCPD с сохранением аннотаций; синтетические сценарии для проверки граничных случаев. \\
Эксперименты & Единые настройки и сохранение сырых результатов; сравнение качества, времени и памяти. \\
Документация & Описание контрактов, пример совместного применения и инструкция воспроизведения. \\
\bottomrule
\end{tabularx}
\end{table}

Исходная цель сохраняется. Выбор первого алгоритма не сужает
назначение обёртки и протокола оценки. Для одного семестра
достижимым результатом является один интегрированный детектор,
работающий оконный прототип, согласованный минимальный набор
метрик и воспроизводимый эксперимент на ограниченном наборе данных.
Дальнейшее расширение множества алгоритмов зависит от результатов
этого этапа.

Предполагаемый вклад работы состоит в расширении возможностей
конкретной библиотеки: согласованных интерфейсах, обработке
последовательного входа и проверяемом протоколе оценки.
Завершённость работы будет определяться совместной проверкой
этих компонентов и возможностью воспроизвести полученные результаты.
